\chapter{Experimental Design and Results}
\label{cha:expdesign}

The training campaign for MMTSFM was still running at the time of writing. This chapter
therefore states, in full, \emph{how} the model is measured --- the arms, the controls, the
reporting rules and the integrity constraints --- and leaves the result tables as
placeholders to be filled from the completed runs. Earlier development runs exist but are
superseded by the current campaign and are deliberately not reported here: a provisional
number that later moves is worse than no number, because it propagates into every
downstream comparison.

Everything below is on the primary scenario of Section~\ref{sec:scenarios}: cross-plant
(S2) on \texttt{uk\_pv}, fourteen disjoint test plants, fourteen days of history, six-hour
horizon, evaluated against the bars established in Section~\ref{sec:bar}.

\section{Arms}
\label{sec:arms}

Four arms constitute the main experiment. They are designed so that each pair differs in
exactly one factor, which is what makes the differences interpretable.

\begin{table}[h!]
\centering
\small
\begin{tabular}{l l l l p{4.4cm}}
\hline
\textbf{Arm} & \textbf{Vision} & \textbf{Mixer} & \textbf{Fusion} & \textbf{Isolates} \\
\hline
A --- numeric & off & Grassmann & --- & The vision-lift lower bound \\
B --- late & on & self-attention & late & The conventional multimodal design \\
C --- interleaved & on & Grassmann & interleaved & The proposed model \\
D --- no pair bias & on & Grassmann & interleaved & The modality-pair correction \\
\hline
\end{tabular}
\caption{The four main arms. All share the curriculum, data, seed and schedule of
Chapter~\ref{cha:method}; they differ only in the stated factor.}
\label{tab:arms}
\end{table}

Three comparisons follow directly:

\begin{itemize}
  \item \textbf{C against A} answers H1 --- does the visual stream contribute measurable
    skill at all? Arm~A is the lower bound: the identical recipe with the vision stack
    disabled. If C does not exceed A, no amount of fusion sophistication is doing work, and
    any advantage over an external baseline comes from the training recipe rather than from
    multimodality.
  \item \textbf{C against B} answers H2 --- does fusion \emph{depth} matter? Same backbones,
    same data, same schedule; only the point at which the modalities meet differs.
  \item \textbf{C against D} isolates the modality-aware offset gating of
    Section~\ref{sec:grassmann}, and is expected to show up in calibration at least as much
    as in point accuracy.
\end{itemize}

Arm~A deserves emphasis because it is the arm most often omitted in multimodal work. Without
a modality-off control run under the \emph{identical} recipe, a multimodal model that beats
an external unimodal baseline proves nothing about its second modality: the gain may come
entirely from the curriculum, the normalisation or the horizon head. Arm~A is what converts
the comparison from suggestive to conclusive, in whichever direction it falls.

\section{Reporting template}
\label{sec:reporting}

Results will be recorded in the form of Table~\ref{tab:mmtsfm-results}, with each arm's
in-process test metrics, and read against the three bars of Section~\ref{sec:bar}: 0.552
overall, 0.540 multimodal, and 0.331 for the same backbone fine-tuned.

\begin{table}[h!]
\centering
\small
\begin{tabular}{l c c c c c}
\hline
\textbf{Arm} & \textbf{NMAE} $\downarrow$ & \textbf{NRMSE} $\downarrow$ &
\textbf{SS} $\uparrow$ & \textbf{CRPS} $\downarrow$ & \textbf{Cov@80} \\
\hline
A --- numeric (vision off) & --- & --- & --- & --- & --- \\
B --- late fusion & --- & --- & --- & --- & --- \\
C --- interleaved (proposed) & --- & --- & --- & --- & --- \\
D --- no pair bias & --- & --- & --- & --- & --- \\
\hline
\multicolumn{6}{l}{\emph{Reference points}} \\
Chronos-2 fine-tuned (same backbone) & 0.1120 & 0.1543 & 0.331 & 0.0855 & --- \\
TS-RAG (best frozen-backbone adaptation) & 0.0705 & 0.1203 & 0.478 & --- & --- \\
iTransformer (overall bar) & 0.0699 & 0.1032 & 0.552 & --- & --- \\
\hline
\end{tabular}
\caption{Reporting template for the MMTSFM campaign. Reference rows are from
Table~\ref{tab:leaderboard}; the arm rows are pending completion of the runs.}
\label{tab:mmtsfm-results}
\end{table}

Three secondary views are reported alongside the aggregate table, because the aggregate
alone cannot distinguish the hypotheses of interest.

\paragraph{Per-horizon profile.} Normalised mean absolute error as a function of lead time,
per arm. If the visual stream contributes, its contribution should be concentrated at short
to medium lead times, where advection is predictable from the observed cloud field, and
should decay as the field leaves the crop. An arm whose profile is indistinguishable from
the vision-free arm at \emph{every} lead time has not used its frames, regardless of what
the aggregate says.

\paragraph{Ramp subset.} The same metrics restricted to the top decile of $|\Delta y|$ per
site (scenario S6). This is where the multimodal hypothesis lives or dies: clear-sky periods
are won by persistence, and vision earns its tokens only during cloud transitions. The
reference to beat here is Time-VLM's ramp NRMSE of 0.172, the best in the suite.

\paragraph{Calibration.} Empirical coverage of the nominal 80\% interval alongside the
continuous ranked probability score. The modality-pair-bias ablation in particular is
expected to affect the predictive distribution more than the conditional mean.

\section{The run matrix}
\label{sec:run-matrix}

Each arm is not one job but a chain of four, one per curriculum stage, each warm-starting from
the previous stage's best checkpoint. Table~\ref{tab:runmatrix} states the full matrix, which
is what determines the compute budget and the order in which results become available.

\begin{table}[h!]
\centering
\small
\begin{tabular}{l c c c c l}
\hline
\textbf{Arm} & \textbf{S1} & \textbf{S2a} & \textbf{S2b} & \textbf{S3} & \textbf{Notes} \\
\hline
A --- numeric (vision off) & yes & --- & --- & yes & S2a/S2b are vision stages and do not
apply \\
B --- late fusion, self-attention & yes & yes & --- & yes & stops at late fusion by
construction \\
C --- interleaved, Grassmann & yes & yes & yes & yes & the full curriculum \\
D --- interleaved, no pair bias & yes & yes & yes & yes & branches from C at S2b \\
\hline
\end{tabular}
\caption{Curriculum stages required per arm. Arms C and D share their S1 and S2a
checkpoints, so the marginal cost of D is two stages rather than four.}
\label{tab:runmatrix}
\end{table}

Two properties of this matrix affect interpretation. Arms C and D share their earlier stages,
so the comparison between them isolates the modality-pair bias cleanly --- they are identical
up to the point where the bias is introduced. Arm A, by contrast, skips the vision stages
entirely and therefore differs from the others in \emph{schedule} as well as in modality; its
S1 and S3 are longer to compensate for the stages it does not run. That asymmetry is
unavoidable --- a vision-free model cannot run a visual alignment stage --- and it is a
caveat on the vision-lift comparison that must be stated whenever that comparison is reported.

\section{Controls}
\label{sec:ablation-design}

Aggregate metrics cannot establish that a model \emph{reads} its second modality; they can
only establish that having it does not hurt. Four controls are specified for that purpose,
and their status is stated honestly.

\begin{table}[h!]
\centering
\small
\begin{tabular}{l p{7.0cm} l}
\hline
\textbf{Control} & \textbf{What it establishes} & \textbf{Status} \\
\hline
Shuffled frames & The frames are read \emph{temporally}. A random permutation of the visual
window should degrade a model that uses advection; a model whose gain is regularisation is
unaffected. & specified \\
Mismatched-plant frames & The frames are \emph{spatially grounded}. Substituting another
plant's frames should degrade a model using local cloud structure; a model relying on a
generic cloudiness prior is unaffected. & specified \\
Modality grid & A four-way decomposition --- target only, target with covariates, target
with vision, all three --- attributing the gain to a specific input rather than to their
combination. & specified \\
Vision-only & An upper bound on the information content of the visual stream alone. &
specified \\
\hline
\end{tabular}
\caption{Control battery. ``Specified'' means designed and registered but not yet executed.}
\label{tab:controls}
\end{table}

The shuffled-frames and modality-grid controls are the highest-value of the four, because
between them they distinguish the three explanations of any observed outcome: the frames
carry usable signal and the model extracts it; the frames carry signal that the encoder
interface fails to transmit; or the frames are redundant with the covariate track, which
already contains cloud cover and irradiance. The benchmark of Chapter~\ref{cha:baselines}
--- where the best multimodal model uses no external imagery --- makes the second and third
explanations live possibilities rather than pedantic alternatives, which raises the value of
running these controls considerably.

Beyond the four above, three sensitivity studies are specified: the visual window length
(three, six and twelve hours), the visual token budget after the latent summariser, and
frozen against partially unfrozen backbones. The last separates the contribution of the
fusion architecture from the contribution of raw adaptation capacity, and complements the
Chronos-2 fine-tuned row already in Table~\ref{tab:leaderboard}.

\section{How each control would be interpreted}
\label{sec:control-logic}

A control is only useful if its possible outcomes have been assigned meanings \emph{before}
it runs. Table~\ref{tab:control-logic} does that for the two decisive controls, so that the
result cannot be rationalised after the fact.

\begin{table}[h!]
\centering
\small
\begin{tabular}{p{3.4cm} p{4.6cm} p{5.4cm}}
\hline
\textbf{Control} & \textbf{Outcome} & \textbf{Conclusion} \\
\hline
Shuffled frames & Skill degrades materially & The model reads temporal structure in the
visual stream; the visual path is working \\
 & Skill unchanged & The frames contribute nothing temporal --- at best a static prior, at
worst nothing \\
 & Skill improves & The visual path is actively harmful and the shuffle acts as
regularisation \\
\hline
Modality grid & Target $<$ target$+$cov $<$ target$+$cov$+$vis & Both modalities contribute;
the architecture works as intended \\
 & target$+$cov $=$ target$+$cov$+$vis & Vision is redundant with the covariate track \\
 & target$+$vis $>$ target$+$cov & Vision carries information the covariates do not, but
fusion of all three fails to combine them \\
\hline
\end{tabular}
\caption{Pre-registered interpretation of the two decisive controls. Assigning meanings in
advance is what prevents a null result from being explained away.}
\label{tab:control-logic}
\end{table}

The third row of each block is the one that would be most informative, and neither is the
outcome the architecture was designed to produce. Evidence from
Chapter~\ref{cha:benchmark} makes them live possibilities rather than pedantic ones: real
satellite frames demonstrably carry cloud-state information beyond the covariate track
(Figure~\ref{fig:cloud-residual}), and yet no model in the suite that consumes them converts
that into skill.

\section{Statistical treatment}
\label{sec:statistics}

Three practices govern how differences between arms are read.

\paragraph{Ties.} Differences smaller than the seed-to-seed variance are reported as ties, not
as orderings. Since that variance has not been measured for this architecture --- one seed per
arm --- the operational rule is conservative: no ordering is claimed between arms differing by
less than approximately 0.01 skill, which is the scale on which nearby leaderboard entries
already sit.

\paragraph{Overlapping windows.} Consecutive forecast origins share most of their context, so
per-window losses are strongly autocorrelated and naive tests overstate significance. The
appropriate instrument is a Diebold--Mariano test~\cite{diebold} on the loss-differential
series with an autocorrelation-consistent variance estimator, computed per plant and combined
across plants. This is specified but not executed, and its absence is a stated limitation.

\paragraph{Per-plant reporting.} Alongside every aggregate figure, the per-plant distribution
is retained, so that an improvement driven by one anomalous plant can be distinguished from a
uniform improvement. Win rate against the reference over the fourteen test plants
(Appendix~\ref{app:aggregation}) is the summary statistic for this: a model beating the
reference on every plant by a modest margin is a more trustworthy result than one with the
same mean driven by three plants.

\section{Compute and reproducibility}
\label{sec:compute}

Training runs on the Leonardo cluster under SLURM, with one job per arm and per-job logs
retained so that any reported number traces to the job that produced it. Configuration is
composed by Hydra; the exact override set for each arm is recorded in the experiment registry
(Appendix~\ref{app:registry}) before the job is submitted, not reconstructed from logs
afterwards.

Three practices support reproducibility beyond seed fixing. Splits are committed to version
control rather than regenerated, and their disjointness is asserted at every load, so a
future run cannot silently receive a different partition. Every experiment is registered
before it launches with a hypothesis, a configuration diff and a nominated comparison
baseline, which prevents the retrospective reframing of an experiment as a test of whatever
it happened to show. And the dataset of record is read-only, so no result can depend on an
undocumented change to the data.

\section{Sensitivity studies}
\label{sec:sensitivity}

Three sweeps are specified beyond the controls. Each varies one design constant that was
fixed by argument rather than by measurement, and each has a prediction attached.

\paragraph{Visual window length: 3, 6 and 12 hours.} Section~\ref{sec:decoupled} argues that
vision carries information only over the cloud-advection horizon, and that widening the
window adds cost rather than signal. The prediction is a concave profile: three hours is too
short to capture an approaching field, twelve hours includes frames whose cloud structure has
left the domain, and six hours --- the configuration of record --- is near the optimum. A
monotone increasing profile would falsify the decoupled-resolution argument and indicate that
the visual context should simply be longer.

\paragraph{Visual token budget.} The latent summariser compresses $4 \times 196$ V-JEPA
patch embeddings into $n_{\text{vis}}$ tokens, with $n_{\text{vis}}=1$ for \texttt{uk\_pv}.
That is aggressive compression, and it is a candidate explanation for a weak visual
contribution: a single 768-dimensional vector may simply be too narrow a channel for the
spatial structure of a cloud field. Sweeping the budget separates ``the vision path does not
work'' from ``the vision path is bottlenecked''. If skill rises with token count, the
interface is the problem; if it is flat, the interface is not.

\paragraph{Frozen against partially unfrozen backbones.} The curriculum unfreezes
progressively, and the final stage trains everything. Comparing arms that stop at each stage
separates the contribution of the fusion architecture from the contribution of raw adaptation
capacity. It also complements the Chronos-2 fine-tuned row of
Table~\ref{tab:leaderboard}: if a fully unfrozen MMTSFM merely matches a fine-tuned bare
backbone, the fusion machinery is contributing nothing that adaptation does not.

\section{Horizon-resolved analysis}
\label{sec:horizon-analysis}

Aggregate metrics average over a six-hour horizon in which the information available to the
model changes character completely. At the first step the model is essentially doing
persistence with a correction; by the twelfth, the cloud field observed at the forecast
origin has largely advected out of the crop and only the numerical context and the
deterministic covariates remain informative. Reporting a single number across that range
hides the mechanism.

Three horizon-resolved views are therefore specified. \emph{Per-step NMAE} per arm, which
should separate the arms at short lead times if vision works at all. \emph{Skill decay} at
one, six and twenty-four hours, which positions the model against the regimes of
Table~\ref{tab:leadtime}. And \emph{the vision-minus-no-vision difference as a function of
lead time}, which is the sharpest single plot the experiment can produce: a positive
difference concentrated at short horizons decaying to zero is exactly the signature the
architecture predicts, and a flat zero line at every horizon is exactly the signature of a
visual path that is not being read.

\section{An efficiency claim measured separately}
\label{sec:efficiency}

The Grassmann operator's asymptotic advantage is a claim about \emph{cost}, and it must be
measured as one rather than inferred from an accuracy table. The comparison against the
matched self-attention twin is therefore reported on two axes: peak memory and step latency
as a function of context length, at $T \in \{128, 512, 1024, 2048\}$ patches, with all else
held fixed.

This matters for how a null accuracy result would be interpreted. If the two mixers score
identically but one is linear in sequence length, the geometric operator is still the better
engineering choice at long context, and the honest conclusion is that its inductive bias is
neutral rather than beneficial at the context lengths tested. Reporting only accuracy would
lose that distinction.

\section{Integrity constraint on reported numbers}
\label{sec:integrity}

One methodological constraint governs every number this chapter will contain, and it is
stated here rather than in a footnote because it restricts what may be claimed.

Saved checkpoints of this model \textbf{do not reproduce their training-time scores when
re-scored in a fresh process}. The discrepancy is large --- validation loss moving from
approximately 2.84 to 4.08, with a corresponding collapse in test skill --- and it affects
both temporal-mixer variants, so it is not specific to the Grassmann operator.

The obvious explanations were eliminated experimentally. Weight restoration is
tensor-faithful: all non-vision parameter keys apply and their values were verified
element-wise after loading. The configurations are bit-identical between the saving and
loading processes. The visual feature cache is unchanged. The dataset iteration is
deterministic under the fixed seed, and the derived buffers are deterministic. Restoring
through the training entry point rather than the evaluation entry point fails identically.
The conclusion that survives is that the checkpoint \emph{content} does not correspond to the
model that produced the scores, despite capturing every registered parameter --- and that
apparently seamless resumed runs are consistent with one or two epochs of retraining
repairing the damage rather than with a faithful restore.

Two consequences follow, and both are binding:

\begin{enumerate}
  \item \textbf{In-process numbers are the record.} Metrics computed within the training
    process that produced the model are internally consistent and are what will be reported.
  \item \textbf{No post-hoc re-scoring.} No number in this thesis may be derived by reloading
    a checkpoint and evaluating it, and no analysis requiring that path --- including
    activation-level inspection of a trained model --- can currently be performed.
\end{enumerate}

The next diagnostic step is an activation fingerprint bisection: dumping per-module outputs
on a fixed batch at save time and again after reload, and locating the first module whose
output diverges. Until that is resolved, the constraint above stands, and the reader should
treat any claim in this thesis that would require re-scoring a checkpoint as unsupported.

\section{What each outcome would mean for the thesis}
\label{sec:outcomes}

The campaign has three broad possible outcomes, and the thesis is written so that each of
them is reportable. Stating this in advance is not hedging; it is what prevents the result
from being framed after the fact to suit whatever happened.

\paragraph{If interleaving beats late fusion, and both beat the vision-free arm.} The
architectural hypotheses H1 and H2 are supported. The contribution is then the mechanism
itself, and the interesting follow-up questions are about scale: does the advantage grow with
context length, where the $O(L)$ mixer's asymptotics start to matter, and does it survive the
cross-dataset shift of scenario S3.

\paragraph{If the arms are indistinguishable.} The vision-free arm matches the fusion arms,
and interleaving matches late fusion. H1 and H2 are unsupported at the scale tested, and the
thesis's contribution shifts to the benchmark and the controls --- which is why
Chapter~\ref{cha:benchmark} is written to stand on its own. This outcome is not a failure of
the experiment; it is a measurement, and given the benchmark evidence that no
satellite-consuming model in the suite converts imagery into skill, it is the outcome the
surrounding evidence currently favours. The controls of Section~\ref{sec:ablation-design}
would then become the primary result, because they distinguish a redundant modality from a
mis-interfaced one.

\paragraph{If the vision-free arm \emph{beats} the fusion arms.} The visual path is actively
harmful --- most likely because the newly initialised summariser and fusion interface consume
capacity and inject noise without supplying information. This would be the most informative
outcome of the three for the field, and the diagnostic path is clear: the token-budget sweep
(Section~\ref{sec:sensitivity}) separates a bottleneck from a nuisance, and the shuffled-frames
control establishes whether anything temporal is being read at all.

In all three cases the reporting rule is the same. The vision-free arm is reported as a
headline row, not as an appendix ablation, and no comparison against an external baseline is
presented without it.

\section{Threats to validity}
\label{sec:threats}

\begin{enumerate}
  \item \textbf{A single dataset.} Every number reported in this thesis is
    \texttt{uk\_pv}. Cross-dataset transfer to \texttt{goes\_pvdaq} (scenario S3) is
    \emph{unproven}, not disproven, and no claim about generalisation beyond a single
    fleet, sensor and climate is supported.
  \item \textbf{A single seed per arm.} Differences smaller than the seed-to-seed variance
    cannot be distinguished, and that variance has not been measured. No significance test
    --- Diebold--Mariano~\cite{diebold} or otherwise --- accompanies the comparisons, so
    small differences between arms must be read as ties.
  \item \textbf{The checkpoint integrity issue} of Section~\ref{sec:integrity}.
  \item \textbf{Ramp subsets are not bit-aligned across tiers.} For the externally
    integrated baselines the ramp mask is recomputed from saved predictions on their native
    windows, so ramp columns are comparable \emph{within} a tier or by rank, not
    element-wise across the whole leaderboard.
  \item \textbf{Two harness-limited baseline rows.} Time-VLM's evaluation windows are not
    aligned with the rest of the suite, and UniCast's weakness is partly attributable to a
    window skew in its native evaluation regime. Both are flagged in
    Section~\ref{sec:implementation}; neither number should be used as a precise comparison
    point.
  \item \textbf{Oracle rows are upper bounds.} The two privileged-information variants
    quantify available headroom and are never competitors.
  \item \textbf{Random rather than region-based plant splits.} As discussed in
    Section~\ref{sec:splits}, spatially adjacent plants may straddle the split boundary,
    which inflates absolute numbers across the board.
\end{enumerate}

\par

